\documentclass[10pt,a4paper]{article}
\usepackage[margin=20mm,headheight=16pt]{geometry}
\usepackage{lmodern,graphicx,booktabs,array,amsmath,xcolor,hyperref,fancyhdr,float,caption,titlesec,microtype}
\renewcommand{\familydefault}{\sfdefault}
\definecolor{ink}{HTML}{233E4A}\definecolor{teal}{HTML}{087F8C}\definecolor{muted}{HTML}{647780}
\color{ink}\hypersetup{colorlinks=true,linkcolor=teal,urlcolor=teal,citecolor=teal}
\setlength{\parindent}{0pt}\setlength{\parskip}{5pt}\setlength{\textfloatsep}{12pt}
\captionsetup{font=small,labelfont={bf,color=teal},labelsep=period,justification=raggedright,singlelinecheck=false,skip=5pt}
\titleformat{\section}{\Large\bfseries\color{teal}}{\thesection}{.6em}{}
\titleformat{\subsection}{\normalsize\bfseries}{\thesubsection}{.5em}{}
\pagestyle{fancy}\fancyhf{}\fancyhead[L]{\small SINGAPORE RAINFALL}\fancyhead[R]{\small Task 1 | 2017--2024}\fancyfoot[L]{\scriptsize All-month analysis | task1-v5}\fancyfoot[R]{\small\thepage}
\renewcommand{\headrulewidth}{.3pt}\graphicspath{{../../figures/task1-v5-report/}}
\newcommand{\plot}[3]{\begin{figure}[H]\centering\includegraphics[width=\linewidth]{#1.pdf}\caption{#2}\label{#3}\end{figure}}
\newcommand{\finding}[1]{\par\smallskip{\color{teal}\bfseries Key observation.} #1\par\smallskip}
\begin{document}
{\color{teal}\small DSA5208 GROUP PROJECT / TASK 1}\par
\vspace{4mm}{\Huge\bfseries Rainfall across Singapore}\par
\vspace{2mm}{\Large Exploring time, place and observational gaps}\par
\vspace{4mm}\textbf{Group members:} [Name 1 / Student ID]; [Name 2 / Student ID, if applicable]; [Name 3 / Student ID, if applicable].

\textbf{Draft prepared:} 9 October 2026. \textbf{Submission deadline:} 18 October 2026.

\section{Research questions and scope}
We analyse the National Environment Agency's five-minute rainfall datasets for 2017--2024 to ask: How does rainfall vary through the year and day? How do comparable observations differ between years? Where do stations record different amounts? The analysis uses Apache Spark for observation-level processing and local Python for visualisation of the resulting small tables.

All eight years and all twelve calendar months are retained in the analytical framework. A year-month without sufficiently complete observations is shown explicitly as unsupported. We do not manufacture missing observations, fill gaps with zero, or label a partial-year total as a complete year. Pair-specific common-station comparisons supplement the descriptive monthly results instead of restricting all analyses to one eight-year station intersection.

\finding{The dataset contains @@OBS@@ cleaned observations from @@STATIONS@@ distinct station IDs. @@ZERO@@\% of recorded observations have zero rain. Data support varies considerably: @@ELIGIBLE@@ of the 96 year-months contain at least one station meeting the 90\% monthly coverage rule. Missingness is therefore part of the interpretation, not just a preprocessing detail.}

\subsection{What each analysis estimates}
\begin{table}[H]\small\centering
\begin{tabular}{p{.22\linewidth}p{.34\linewidth}p{.34\linewidth}}\toprule Question & Sample and quantity & Interpretation limit\\\midrule
Seasonal variation & Qualified station-months; mean monthly station totals, wet frequency and intensity & Stations and available years vary\\
Calendar-year totals & Annual coverage $\geq95\%$ and each month $\geq90\%$; station totals and their mean & Unsupported years stay blank; no national area weighting\\
Cross-year differences & Common stations for a particular year pair and calendar month & Different pairs need not contain the same stations\\
Spatial variation & Qualified stations in each of 96 year-months; point maps and planning-area station means & Unsupported locations are not interpolated or assigned zero\\\bottomrule
\end{tabular}\caption{Analysis questions are separated so that an available-sample summary is not mistaken for a fully balanced comparison.}\end{table}

This document covers Task 1 only. Forecast model features, training and predictive evaluation belong to the subsequent Task 2 report. Group details and the final interpretation require author review before submission.

\clearpage
\section{Data exploration}
\subsection{Scale, measurement and rainfall distribution}
Each record identifies a station, observation time, coordinates and rainfall amount in millimetres, together with update metadata. Times are interpreted in Asia/Singapore time. A recorded zero indicates no measured rain during that observation interval; absence of a record does not imply dry weather. The publisher notes that records can be missing and have not received the quality control used for official climate records \cite{data}.
\plot{01_data_profile}{Observation counts, the proportion of recorded zeros, and quantiles of positive five-minute rainfall. Each panel answers a different question about the sample; the right panel excludes zero observations.}{fig:profile}
\finding{Rainfall is strongly concentrated at zero, while the positive measurements have a long upper tail. This motivates examining occurrence frequency separately from the amount conditional on rain, rather than relying on one mean alone.}
\plot{02_positive_distribution}{Distribution of positive five-minute amounts over all eight years. The horizontal axis is logarithmic and bars show the fraction of positive observations in each log-spaced bin. Zero rainfall is summarised separately in Figure~\ref{fig:profile}.}{fig:hist}
The largest retained reading is @@MAX@@ mm in five minutes. High values are flagged for inspection, not automatically removed by an arbitrary cutoff. Neither an empirical percentile nor a neighbouring positive observation independently verifies a meteorological extreme. Five-minute amounts must not be confused with hourly rates or daily accumulations.

\clearpage
\subsection{Coverage and the observation calendar}
For station $s$ in year $y$ and month $m$, coverage is
\[
 C_{sym}=\frac{N_{sym}}{288\times D_{ym}},\qquad 288=\frac{24\times60}{5}.
\]
$N_{sym}$ counts valid unique observations and $D_{ym}$ is the number of calendar days. A calendar skeleton includes absent station-months, with zero observed counts but undefined rainfall totals. This also means partial station lifetimes can lower calendar coverage without implying equipment failure throughout every absent period.
\plot{03_coverage}{Monthly coverage and analysis support. The upper panel is the median among stations with any records. The lower panel counts stations reaching 90\% coverage; grey indicates none qualify. Station counts and completeness must be read together.}{fig:coverage}
\plot{04_source_gap}{Daily raw-versus-clean counts for August 2018. Absence after 3 August already occurs in the source files; the cleaning pipeline did not create this gap.}{fig:gap}
The audit finds @@EMPTY@@ calendar dates with no parsed raw observations anywhere in the network. This identifies source-file gaps but not their operational cause. For August 2018, @@AUGRAW@@ raw records occur on only the first three days. Such a month cannot support an observed full-month total comparable with well-covered months.

\clearpage
\section{Cleaning and analytical methods}
\subsection{Auditable cleaning}
The raw annual files are preserved. PySpark standardises station identifiers, parses timestamps, checks measurement type and units, and casts rainfall and coordinates to numeric values. Malformed records, invalid timestamps, non-finite or negative rainfall and invalid coordinates are handled explicitly. Timestamps within two seconds of a five-minute boundary are snapped to that boundary; off-grid records outside that tolerance are excluded.

Exact duplicates are removed. Conflicting rainfall or coordinate measurements at the same station-time key are quarantined rather than arbitrarily selected. Consistent measurements with repeated update metadata are consolidated; conservative publication timestamps are retained. Large non-negative rainfall is retained for review. Clean data are saved as Parquet partitioned by year and month, accompanied by quality and inspection outputs.

The fresh daily audit reconciles @@RAW@@ source rows with @@OBS@@ cleaned observations, a difference of @@REMOVED@@ rows; this difference includes consolidation and must not all be described as invalid measurements. Raw timestamps that could not be parsed in this audit: @@PARSE@@. Raw daily counts use source timestamps; cleaning can move a few near-midnight observations into an adjacent date when snapping to the five-minute grid, so daily differences are not all deletions. Independent checks reconcile daily, monthly and yearly counts and reproduce aggregate and pairwise statistics from the saved tables.

\subsection{Exposure, frequency and intensity}
Monthly summaries use station-months with coverage $\geq90\%$. This is a predeclared completeness rule, not a guarantee of random missingness. Observed monthly station totals are averaged equally across eligible stations. Wet frequency and conditional intensity use valid observation counts:
\[
 f=\frac{N_{\rm wet}}{N_{\rm observed}},\qquad
 I=\frac{\sum_i r_i}{N_{\rm wet}},\qquad \bar r=fI.
\]
A wet observation has $r_i>0$. Equal-year means are used for the seasonal monthly curves; diurnal summaries pool observed intervals. These weighting schemes are stated because they estimate different summaries. No total is scaled by inverse coverage to manufacture unobserved rainfall.

For cross-year comparisons, eligible observations are matched by station and calendar month within each year pair. Station-month coordinate spans and between-period mean coordinates are screened at 0.005 degrees. The paired change is the average of within-station changes, with percentage change measured relative to the earlier pair mean. Plots suppress changes supported by fewer than five stations; this is a display rule, not a statistical significance threshold. All pair counts remain in the output. Annual comparisons require the full-year eligibility rule in both years.

\subsection{Computing environment and reproducibility}
\begin{table}[H]\small\centering\begin{tabular}{p{.29\linewidth}p{.62\linewidth}}\toprule Component & Recorded configuration\\\midrule
Compute & Google Cloud Dataproc, asia-southeast1-c; one master and two workers\\
VMs and disks & Each e2-standard-4: 4 vCPUs, 16 GB RAM, 80 GB pd-balanced boot disk\\
Worker capacity & 8 vCPUs and 32 GB RAM; master resources are separate\\
Runtime & Spark 3.5.3, Dataproc 2.2.87 Debian 12, PySpark on YARN\\
Spark settings & 6 GB driver memory, 1 GB maximum driver result, 128 shuffle partitions, adaptive SQL, Asia/Singapore timezone\\
Storage and graphics & Cloud Storage raw/Parquet/result tables; local Apple M2, 16 GB RAM; pandas, Matplotlib and Plotly\\\bottomrule\end{tabular}
\caption{The Spark job performs record-level work. Only aggregated tables are downloaded for plotting.}\end{table}
Run-specific source snapshots, metadata and validation outputs accompany \texttt{task1-v5}. Figures are generated as vector PDF/SVG and PNG previews. Report captions and references are typeset in LaTeX rather than embedded as miniature text in chart images.

\clearpage
\section{Monthly and seasonal rainfall patterns}
\plot{05_monthly_amount}{Monthly mean observed station totals for all 2017--2024 year-months. A common linear colour scale supports magnitude comparisons. Grey cells lack qualified station-months and must not be read as dry weather.}{fig:monthly}
\finding{Among eligible-year monthly averages, @@WETMONTH@@ has the largest mean amount (@@WETMM@@ mm) and @@DRYMONTH@@ the smallest (@@DRYMM@@ mm). Each year-month first averages its qualified stations; the seasonal average then weights available years equally.}
\plot{06_seasonality}{Monthly amount, wet-interval frequency and conditional wet intensity. The shaded region is the observed minimum--maximum across eligible years, not a confidence interval.}{fig:seasonal}
The seasonal distinction is not solely about intensity: a month can be wet because rain occurs more frequently, because individual wet intervals contain more rain, or both. MSS describes generally wetter conditions in November--January and a relatively dry February \cite{climate}. The present sample is a short, incomplete observational period, not an official climate normal. Calendar months also differ in length.

\clearpage
\subsection{The same month across eight years}
\plot{07_monthly_trends}{Separate trajectories for every calendar month. Missing year-months break the lines. Vertical ranges differ between panels to reveal within-month variation; use Figure~\ref{fig:monthly} for comparisons of absolute magnitude between months.}{fig:trajectories}
This view preserves usable observations separately for each calendar month. For example, an unsupported August 2018 does not remove the other seven Augusts, nor does it remove well-covered months in 2018. However, individual points can involve different station compositions, so changes in these descriptive curves are not automatically changes at fixed locations.
\plot{09_paired_changes}{Adjacent-year monthly changes using common stable-location stations. Values are percentage changes in the pair's mean observed station amount. Grey denotes missing support or fewer than five common stations; the colour scale is centred on zero.}{fig:paired}
@@PAIRTEXT@@ The pairing controls station identity for each comparison. It does not control unobserved rainfall within remaining gaps, and pair-specific panels differ. Consequently, these percentages are not chained to construct a national trend, and no significance claim is attached to them.

\clearpage
\section{Calendar-year rainfall totals}
\plot{08_annual_totals}{Observed full-calendar-year amounts at eligible stations. Dots show individual stations and coral diamonds their mean; $n$ is the number of stations. Shaded columns mark years with no station meeting the annual and every-month completeness rules.}{fig:annual}
\begin{table}[H]\centering\small\begin{tabular}{rrrr}\toprule Year & Qualified stations & Mean total (mm) & Mean coverage (\%)\\\midrule
@@ANNUALROWS@@
\bottomrule\end{tabular}\caption{Annual summaries describe the eligible network, not a sum over all stations or an area-weighted national measurement. Unsupported years remain in the calendar but have no qualifying mean.}\end{table}
@@ANNUALTEXT@@

A calendar-year total is the sum of a station's valid five-minute amounts over January--December. Even eligible totals are observed amounts and retain residual gaps. Combining station totals by addition would make the result depend strongly on the number of gauges; the plotted network statistic is therefore a station mean. Spatially changing station sets are another reason not to infer a national wettest-year ranking from these values.

The pairwise annual output also compares common eligible stations between supported years. @@ANNUALPAIR@@ This is a restricted geographical comparison and does not recover the missing annual totals in other years. Annual coverage alone is insufficient: a station can appear well covered overall while missing a concentrated part of a season, which motivates the additional monthly requirement.

\clearpage
\section{Diurnal behaviour and heavier daily rainfall}
\plot{10_diurnal}{Wet-observation frequency by month and hour (left) and year and hour (right), in Singapore time. All eligible station-months contribute; values are weighted by observed intervals. The year curves can include different month mixes.}{fig:diurnal}
@@HOURTEXT@@ This is descriptive occurrence frequency, not a probability forecast for a future event. The source's incomplete time coverage and changing station mix remain relevant even after screening. The general afternoon enhancement is consistent with MSS's description of Singapore's daily rainfall pattern \cite{climate}; it does not establish the cause of a particular storm.
\plot{11_daily_tails}{Monthly daily-rainfall P99 (upper) and wet station-day percentage (lower). Only days with at least 95\% coverage inside eligible station-months are included. Quantiles include dry days; wet means an observed daily amount above zero.}{fig:tails}
High quantiles provide a data-derived description of the upper tail without declaring an arbitrary amount to be an official extreme threshold. They are approximate Spark quantiles of station-days, not independent storms or return levels. A high wet-day proportion and a high upper-tail amount need not coincide.

\clearpage
\section{Spatial patterns across the eight years}
All 96 year-month combinations have station-map views. Figures~\ref{fig:spatialdec} and \ref{fig:regions} use December as a like-for-like calendar illustration; the gallery provides corresponding comparisons for every other month. All monthly maps use one linear scale over the full 96-period set, preserving comparability and retaining high values even when this reduces contrast among moderate observations.
\plot{13_spatial_month_12}{December station observations across all eight years. Labels show qualified station counts. Colours represent observed monthly totals; the maps contain no interpolated rainfall.}{fig:spatialdec}
\plot{15_regions_month_12}{December planning-area heatmaps across eight years. Each coloured polygon displays the arithmetic mean of assigned eligible stations; small dots show contributing locations. Grey means no supporting station, not zero rainfall.}{fig:regions}
@@SPATIALTEXT@@ Uneven station placement and different month-specific networks preclude treating every coloured polygon as an equally reliable area-average measurement. Per-area station counts and assignments are retained in the supplementary CSVs; boundary or ambiguous points are excluded from polygon means rather than forced into a region.

\clearpage
\subsection{Spatial seasonality and annual support}
\plot{14_spatial_year_2024}{All twelve months of 2024, the final study year. The same global monthly colour scale is used as in the eight-year comparisons. Other years have equivalent twelve-panel figures in the gallery.}{fig:spatialseason}
Month-specific eligibility makes the twelve-month view possible without requiring every station to survive an eight-year intersection. Within a year, the qualified station set can still change between months; the displayed counts expose this. The complete interactive station map allows inspection of rainfall and coverage at individual sites for every supported period.
\plot{16_annual_spatial}{Annual observations using the stricter full-year completeness rule. Unlike the monthly maps, several annual panels have no qualifying stations. Blank panels communicate the limit of annual evidence rather than restricting all spatial exploration to one selected year.}{fig:annualmaps}
Boundary geometries are derived from the retained geoBoundaries/URA assets. Their vintage is older than parts of the observation period; coastal alignment and historical administrative definitions can affect assignments. We prioritise station observations and explicitly supported regional summaries over smooth interpolation that could imply more spatial information than was measured.

\clearpage
\section{Sensitivity, limitations and additional observations}
\plot{12_sensitivity}{Seasonal amount and wet-frequency summaries repeated at three monthly coverage cutoffs. These compare rainfall results, not only the number of surviving samples. Each curve still uses its own eligible stations and years.}{fig:sensitivity}
@@SENSTEXT@@ Comparing matched year-months isolates changes from the availability of whole months, but the contributing stations within them still change. Similar curves are evidence of stability for these particular descriptive summaries, not proof that missingness is ignorable or that all spatial and paired conclusions are robust.

The principal limitations are: (i) source gaps may coincide with particular weather and are not known to be random; (ii) eight years are too short and incomplete here to establish a climate-change trend; (iii) station means are not area-weighted national rainfall; (iv) empirical tails can contain measurement or export errors; and (v) pairs and regional summaries have varying geographical support. High observations require review against independent sources before being presented as verified extreme events. Future work could compare an explicitly justified missing-data model against complete-period evidence, with uncertainty reported rather than concealed.

For Task 2, the high zero proportion motivates evaluation beyond accuracy, while occurrence seasonality and station history suggest candidate predictors. Training eligibility should use complete historical and future windows around each prediction time, not the strict eight-year station intersection. No forecast model performance is claimed in this Task 1 report.

\section{AI usage declaration}
OpenAI Codex assisted with code development and revision, data-quality diagnostics, Cloud execution, figure design and this draft. Numerical statements were derived from saved outputs; AI-generated interpretation does not independently validate measurements. The group must review the source code, results and references, fill in the author details, and take responsibility for the submitted work.

\begin{thebibliography}{9}
\bibitem{data} National Environment Agency. \emph{Historical Rainfall across Singapore}, annual datasets 2017--2024. Data.gov.sg. \url{https://data.gov.sg/collections/2279/view}. Accessed 9 October 2026.
\bibitem{climate} Meteorological Service Singapore. \emph{Climate of Singapore}. \url{https://www.weather.gov.sg/climate-climate-of-singapore/}. Accessed 9 October 2026.
\bibitem{geo} geoBoundaries / Urban Redevelopment Authority. Singapore outline and 55 planning-area geometries. Provenance and licences retained in \texttt{assets/singapore.source.json} and \texttt{assets/planning\_areas.source.json}; planning-area source is based on URA Master Plan 2014 geometry.
\end{thebibliography}
\end{document}
